%--------------------------------------------------------------------%
% Page          : Abstract (English)
%--------------------------------------------------------------------%

\clearpage

\addcontentsline{toc}{chapter}{ABSTRACT}

\section*{\large \centering \MakeUppercase 
    {Abstract}
}

\begin{center}

    \justifying \normalsize

    \textbf{Decision Making and Elastic Scalability Management in Heterogeneous Cloud Environments Based on Workload Prediction}

    \vspace{0.5cm}

    \textbf{Muhammad Zahid Masruri / NRP 6025222041}

    \vspace{0.5cm}

    Modern cloud architecture faces a trade-off between scalability and cost efficiency. Kubernetes provides strong infrastructure control but is slow to respond to traffic spikes, while serverless computing offers high elasticity at a higher per-request cost. This research designs, implements, and evaluates a hybrid architecture integrating Kubernetes and serverless clusters with Gated Recurrent Unit (GRU)-based workload prediction for proactive traffic distribution.

    The proposed system implements a V3 capacity-driven routing controller that dynamically shifts traffic between Kubernetes and serverless backends based on tail latency monitoring. The proactive routing mechanism uses actual workload trend extrapolation (5 steps $\times$ 15 seconds = 75 seconds ahead) gated by GRU confidence. Evaluation is conducted through 30 replicated controlled experiments across four scenarios: K8s+HPA (baseline), Serverless-only, Hybrid-reactive, and Hybrid-predictive, using the ClarkNet HTTP trace (40 stages $\times$ 30 seconds = 20 minutes, RPS 22--164).

    Results show S4 (Hybrid-predictive) achieves p99 latency = 188\,ms with 702 SLO violations per run, 75\% fewer than S3 (Hybrid-reactive) with 2,805 violations (Cohen's d = 1.21, large effect). S4 is also 20\% cheaper than S1 (\$149/month vs \$187/month) at 98.4\% SLO compliance. The GRU model was optimized via Optuna TPE HPO (30 trials, RMSE 6.01\% $\rightarrow$ 4.75\%) achieving 400/400 successful predictions with inference latency 10--13\,ms. Methodological finding: single-run HPO on the online controller caused overfitting detectable only through $n=5$ holdout validation, demonstrating the necessity of capacity-test-based calibration and multi-run confirmation for shared-resource systems.

    \vspace{0.5cm}

    \textbf{Keywords:} \emph{Cloud Computing, Kubernetes, Serverless, Workload Prediction, GRU, Traffic Distribution, Hybrid Architecture}

\end{center}

\clearpage
